%=====================================================================
%  Quantitative Reasoning - LAB 413
%  Lab 9: Inference for Proportions & Difference in Two Proportions
%  Instructor: Subin Na
%  Beamer conversion of Lab9_1105_PPT.pptx
%  Compile: pdflatex Lab9_beamer.tex   (run twice for footline counters)
%=====================================================================
\documentclass[aspectratio=169,11pt]{beamer}

\usepackage[utf8]{inputenc}
\usepackage[T1]{fontenc}
\usepackage{lmodern}
\usepackage{amsmath,amssymb}
\usepackage{graphicx}
\usepackage{booktabs}
\usepackage{tcolorbox}
\tcbuselibrary{skins}

\usetheme{default}
\usecolortheme{default}
\useinnertheme{circles}

\definecolor{qrnavy}{RGB}{20,44,90}
\definecolor{qrblue}{RGB}{38,80,150}
\definecolor{qrgray}{RGB}{90,95,105}
\definecolor{qrlight}{RGB}{234,239,247}
\definecolor{qrred}{RGB}{190,60,60}

\setbeamercolor{structure}{fg=qrnavy}
\setbeamercolor{frametitle}{fg=white,bg=qrnavy}
\setbeamercolor{title}{fg=qrnavy}
\setbeamercolor{block title}{fg=white,bg=qrblue}
\setbeamercolor{block body}{bg=qrlight}
\setbeamercolor{itemize item}{fg=qrblue}
\setbeamercolor{itemize subitem}{fg=qrblue}
\setbeamerfont{frametitle}{series=\bfseries}
\setbeamertemplate{navigation symbols}{}

\setbeamertemplate{frametitle}{%
  \nointerlineskip
  \begin{beamercolorbox}[wd=\paperwidth,ht=3.2ex,dp=1.2ex,leftskip=0.3cm]{frametitle}%
    \strut\insertframetitle\strut
  \end{beamercolorbox}%
}

% Footline: course | lab N | slide number  (edit "Lab N" per deck)
\newcommand{\labnum}{Lab 9}
\setbeamertemplate{footline}{%
  \leavevmode%
  \hbox{%
    \begin{beamercolorbox}[wd=.5\paperwidth,ht=2.4ex,dp=1ex,leftskip=.3cm]{author in head/foot}%
      \usebeamerfont{author in head/foot}\color{qrgray}\small Quantitative Reasoning -- LAB 413%
    \end{beamercolorbox}%
    \begin{beamercolorbox}[wd=.5\paperwidth,ht=2.4ex,dp=1ex,rightskip=.3cm]{title in head/foot}%
      \usebeamerfont{title in head/foot}\color{qrgray}\small\hfill \labnum \quad\insertframenumber\,/\,\inserttotalframenumber%
    \end{beamercolorbox}}%
}

% Number badge (01, 02, ...) used on concept slides:  \numbadge{01}
\newtcbox{\numbadge}{on line,boxsep=2pt,left=4pt,right=4pt,top=2pt,bottom=2pt,
  colback=qrnavy,colframe=qrnavy,fontupper=\bfseries\color{white},arc=2pt}

% Key-concepts callout:  \begin{keybox} ... \end{keybox}   or  \begin{keybox}[Scenario] ...
\newtcolorbox{keybox}[1][Key Concepts]{colback=qrlight,colframe=qrnavy,
  fonttitle=\bfseries,title={#1},coltitle=white,arc=2pt,left=4pt,right=4pt,
  top=3pt,bottom=3pt,boxrule=0.6pt}

%---------------------------------------------------------------------
%  Title data
%---------------------------------------------------------------------
\title{\bfseries Inference for Proportions \& Difference in Two Proportions}
\subtitle{Week 11 Recap \textbullet\ Lab 9}
\author{Instructor: Subin Na}
\institute{Quantitative Reasoning -- LAB 413}
\date{November 5, 2025}

\begin{document}

%---------------------------------------------------------------------
%  Title slide
%---------------------------------------------------------------------
{\setbeamertemplate{footline}{}
\begin{frame}
  \vfill
  {\color{qrgray}\small Quantitative Reasoning -- LAB 413}\\[0.6em]
  {\usebeamerfont{title}\color{qrnavy}\huge\bfseries Inference for Proportions\\[0.15em]\& Difference in Two Proportions\par}
  \vspace{0.4em}
  {\color{qrblue}\large Week 11 Recap \ \textbullet\ \ Lab 9}\par
  \vspace{1.6em}
  {\color{qrgray} Instructor: \textbf{Subin Na} \hfill November 5, 2025}
  \vfill
\end{frame}
}

%---------------------------------------------------------------------
%  Today's Plan
%---------------------------------------------------------------------
\begin{frame}{Today's Plan}
  \begin{columns}[T]
    \begin{column}{0.5\textwidth}
      \textbf{\color{qrnavy}1. Quick Recap}
      \begin{itemize}
        \item Bernoulli \& Binomial Foundations
        \item Sampling Distribution of $\hat{p}$
        \item Confidence Interval for a Single Proportion
        \item Hypothesis Test for a Single Proportion
        \item Comparing Two Proportions
      \end{itemize}
    \end{column}
    \begin{column}{0.5\textwidth}
      \phantom{\textbf{x}}
      \vspace{4.2em}

      \textbf{\color{qrnavy}2. RStudio Practice}
    \end{column}
  \end{columns}
\end{frame}

%---------------------------------------------------------------------
%  01 Bernoulli & Binomial Foundations
%---------------------------------------------------------------------
\begin{frame}{\numbadge{01}\ \ Bernoulli \& Binomial Foundations}
  \begin{itemize}
    \item \textbf{Bernoulli trial:} one trial $\rightarrow$ outcome $0$ or $1$ (failure/success).
    \item \textbf{Binomial distribution:} $X = $ number of successes in $n$ independent Bernoulli trials.
  \end{itemize}
  \begin{equation*}
    E[X] = np, \qquad \mathrm{Var}[X] = np(1-p)
  \end{equation*}
  \begin{itemize}
    \item \textbf{Sample proportion} $\hat{p} = X/n$ $\rightarrow$ point estimate for the population proportion $p$.
  \end{itemize}
  \begin{keybox}[Example]
    Among $n = 200$ voters, $X = 120$ support Candidate A
    $\;\Rightarrow\; \hat{p} = 120/200 = 0.60$.
  \end{keybox}
\end{frame}

%---------------------------------------------------------------------
%  02 Sampling Distribution of p-hat
%---------------------------------------------------------------------
\begin{frame}{\numbadge{02}\ \ Sampling Distribution of $\hat{p}$}
  \begin{keybox}
    If the sample is random and large enough ($n\hat{p} \ge 10$ and $n(1-\hat{p}) \ge 10$),
    the sampling distribution of $\hat{p}$ is approximately Normal.
  \end{keybox}
  \vspace{0.6em}
  \begin{equation*}
    \hat{p} \;\sim\; N\!\left(p,\ \sqrt{\frac{p(1-p)}{n}}\right)
  \end{equation*}
  \begin{itemize}
    \item $\hat{p}$ is an \textbf{unbiased estimator}: on average $\hat{p} = p$.
    \item The \textbf{standard error} $\sqrt{p(1-p)/n}$ shrinks as $n$ grows $\rightarrow$ more stable estimates.
  \end{itemize}
\end{frame}

%---------------------------------------------------------------------
%  03 Confidence Interval for a Single Proportion
%---------------------------------------------------------------------
\begin{frame}{\numbadge{03}\ \ Confidence Interval for a Single Proportion}
  \begin{columns}[T]
    \begin{column}{0.52\textwidth}
      \begin{equation*}
        \hat{p} \;\pm\; z^{*}\sqrt{\frac{\hat{p}(1-\hat{p})}{n}}
      \end{equation*}
      {\small where $z^{*} \approx 1.645$ (90\%), $1.96$ (95\%), $2.576$ (99\%).}
      \vspace{0.5em}
      \begin{keybox}[Plus-Four Adjustment]
        \small If counts are small, add 2 successes and 2 failures:
        $\displaystyle \tilde{p} = \frac{X+2}{n+4}$.
      \end{keybox}
    \end{column}
    \begin{column}{0.46\textwidth}
      \textbf{\color{qrnavy}Conditions}
      \begin{itemize}
        \item Random sample.
        \item At least 10 expected successes \& failures.
      \end{itemize}
      \vspace{0.4em}
      \textbf{\color{qrnavy}Interpretation}
      \begin{itemize}
        \item If we repeatedly sampled, $\approx C\%$ of intervals would capture the true $p$.
      \end{itemize}
    \end{column}
  \end{columns}
\end{frame}

%---------------------------------------------------------------------
%  Example: Dogs vs Cats - setup
%---------------------------------------------------------------------
\begin{frame}{Example: Dogs vs.\ Cats --- Setup}
  \begin{keybox}[Scenario]
    A quick campus poll asks: ``Do you prefer dogs over cats?''
    \begin{itemize}
      \item Random sample of $n = 600$ students.
      \item $x = 372$ answered ``Yes.''
    \end{itemize}
  \end{keybox}
  \vspace{0.6em}
  \begin{center}
    \textbf{Goal:} estimate the true proportion $p$ of all students who prefer dogs,
    and construct a \textbf{90\% confidence interval} for $p$.
  \end{center}
\end{frame}

%---------------------------------------------------------------------
%  Example: Dogs vs Cats - solution
%---------------------------------------------------------------------
\begin{frame}{Example: Dogs vs.\ Cats --- Solution}
  \small
  \textbf{\color{qrnavy}Given.} $n = 600,\ x = 372,\ \hat{p} = 372/600 = 0.62$.

  \textbf{\color{qrnavy}Conditions.} $n\hat{p} = 372 > 10,\ n(1-\hat{p}) = 228 > 10.$ \checkmark
  \vspace{0.3em}
  \begin{align*}
    SE &= \sqrt{\frac{\hat{p}(1-\hat{p})}{n}} = \sqrt{\frac{0.62(0.38)}{600}} = 0.0196\\[0.3em]
    \hat{p} \pm z^{*} SE &= 0.62 \pm 1.645(0.0196) = 0.62 \pm 0.032 \;\Rightarrow\; (0.588,\ 0.652)
  \end{align*}
  \begin{tcolorbox}[colback=qrlight,colframe=qrnavy,arc=2pt,boxrule=0.6pt]
    \small The 90\% CI for the proportion who prefer dogs over cats is $(58.8\%,\ 65.2\%)$,
    based on a method that would capture the true proportion about 90\% of the time.
  \end{tcolorbox}
\end{frame}

%---------------------------------------------------------------------
%  04 Hypothesis Test for a Single Proportion
%---------------------------------------------------------------------
\begin{frame}{\numbadge{04}\ \ Hypothesis Test for a Single Proportion}
  \begin{columns}[T]
    \begin{column}{0.5\textwidth}
      \begin{block}{Hypotheses}
        \[ H_0: p = p_0 \qquad H_a: p \neq p_0 \]
      \end{block}
    \end{column}
    \begin{column}{0.5\textwidth}
      \begin{block}{Test statistic}
        \[ z = \frac{\hat{p} - p_0}{\sqrt{p_0(1-p_0)/n}} \]
      \end{block}
    \end{column}
  \end{columns}
  \vspace{0.4em}
  \begin{keybox}
    Reject $H_0$ when the test statistic is more extreme than the critical value
    --- equivalently, when the $p$-value is smaller than $\alpha$.
  \end{keybox}
\end{frame}

%---------------------------------------------------------------------
%  04 Example: Boba tea vs coffee
%---------------------------------------------------------------------
\begin{frame}{Example: Boba Tea vs.\ Coffee}
  \small
  \textbf{\color{qrnavy}Setup.} Out of $n = 20$ students, $13$ prefer boba tea over coffee,
  so $\hat{p} = 0.65$. Test whether this differs from $0.5$.
  \vspace{0.3em}
  \begin{align*}
    z &= \frac{0.65 - 0.5}{\sqrt{0.5(1-0.5)/20}} = 1.34\\[0.3em]
    p\text{-value} &= 2\,P(Z > 1.34) = 2 \times 0.0901 = 0.18
  \end{align*}
  \begin{tcolorbox}[colback=qrlight,colframe=qrred,arc=2pt,boxrule=0.6pt]
    \small $p = 0.18 > \alpha = 0.05 \;\Rightarrow\;$ \textbf{fail to reject} $H_0$.
    Not enough evidence that the true proportion differs from $0.5$.
    (``Fail to reject'' does \emph{not} prove $H_0$ true.)
  \end{tcolorbox}
\end{frame}

%---------------------------------------------------------------------
%  05 Comparing Two Proportions
%---------------------------------------------------------------------
\begin{frame}{\numbadge{05}\ \ Comparing Two Proportions}
  \textbf{Goal:} test or estimate the difference between two independent populations.
  \begin{equation*}
    (\hat{p}_1 - \hat{p}_2) \;\pm\; z^{*}\sqrt{\frac{\hat{p}_1(1-\hat{p}_1)}{n_1} + \frac{\hat{p}_2(1-\hat{p}_2)}{n_2}}
  \end{equation*}
  \begin{columns}[T]
    \begin{column}{0.5\textwidth}
      \textbf{\color{qrnavy}Conditions}
      \begin{itemize}
        \item Two independent random samples.
        \item Each sample has $\ge 10$ expected successes \& failures.
      \end{itemize}
    \end{column}
    \begin{column}{0.5\textwidth}
      \textbf{\color{qrnavy}Interpretation}
      \begin{itemize}
        \item CI includes 0 $\rightarrow$ no significant difference.
        \item CI excludes 0 $\rightarrow$ significant difference.
      \end{itemize}
    \end{column}
  \end{columns}
\end{frame}

%---------------------------------------------------------------------
%  Example: Pineapple on pizza - setup
%---------------------------------------------------------------------
\begin{frame}{Example: Pineapple on Pizza --- Setup}
  \begin{keybox}[Scenario]
    Among students, does support for pineapple on pizza differ by gender?
    Two independent samples:
    \begin{itemize}
      \item Women: $328$ of $537$ said ``yes.''
      \item Men: $234$ of $532$ said ``yes.''
    \end{itemize}
  \end{keybox}
  \vspace{0.6em}
  \begin{center}
    \textbf{Goal:} construct a \textbf{95\% confidence interval} for the difference $p_1 - p_2$
    and decide whether the difference is statistically significant.
  \end{center}
\end{frame}

%---------------------------------------------------------------------
%  Example: Pineapple on pizza - solution
%---------------------------------------------------------------------
\begin{frame}{Example: Pineapple on Pizza --- Solution}
  \small
  \textbf{\color{qrnavy}Given.} $\hat{p}_1 = 0.611$ (women), $\hat{p}_2 = 0.440$ (men).
  Both samples meet the $\ge 10$ successes \& failures condition. \checkmark
  \begin{align*}
    SE &= \sqrt{\frac{0.611(1-0.611)}{537} + \frac{0.440(1-0.440)}{532}} = 0.030\\[0.3em]
    (\hat{p}_1 - \hat{p}_2) \pm z^{*} SE &= (0.611 - 0.440) \pm 1.96(0.030) = 0.171 \pm 0.059\\[0.2em]
    &\Rightarrow\; 95\%\ \text{CI} = (0.112,\ 0.230)
  \end{align*}
  \begin{tcolorbox}[colback=qrlight,colframe=qrnavy,arc=2pt,boxrule=0.6pt]
    \small The 95\% CI for the difference (women $-$ men) is $(0.112,\ 0.230)$.
    Since it excludes 0, women are significantly more likely to approve pineapple on pizza.
  \end{tcolorbox}
\end{frame}

%---------------------------------------------------------------------
%  06 Key Takeaways
%---------------------------------------------------------------------
\begin{frame}{\numbadge{06}\ \ Key Takeaways}
  \begin{keybox}
    \begin{itemize}
      \item When $n$ is large enough, the sampling distribution of $\hat{p}$ is approximately Normal $\rightarrow$ we can use $z$ procedures.
      \item A confidence interval is always \textbf{estimate $\pm$ margin of error}: for proportions, $\hat{p} \pm z^{*} \times SE$.
      \item A hypothesis test compares the sample statistic to the null value using a $z$-score (distance in standard-error units).
      \item Always check conditions first: random sample and large-sample thresholds ($\ge 10$ expected successes \& failures).
      \item Two-proportion procedures follow the same logic as two-sample mean tests --- only the formulas change.
    \end{itemize}
  \end{keybox}
\end{frame}

%---------------------------------------------------------------------
%  Transition to RStudio
%---------------------------------------------------------------------
{\setbeamertemplate{footline}{}
\begin{frame}
  \vfill
  \centering
  {\color{qrnavy}\Huge\bfseries RStudio Practice}\\[0.8em]
  {\color{qrblue}\large Comparing Means: Two-Sample $t$-test (McGrath Experiment)}\\[0.4em]
  {\color{qrgray} See \texttt{Lab9.Rmd} / \texttt{Lab9\_1105.R}}
  \vfill
\end{frame}
}

\end{document}
